\RequirePackage{iftex}
\ifPDFTeX
\documentclass[a4paper,12pt]{article}
\usepackage[hmargin=16mm,vmargin=24mm]{geometry}
\usepackage[T1]{fontenc}
\usepackage[utf8]{inputenc}
\usepackage{CJKutf8}
\else
\documentclass[a4paper,12pt]{jsarticle}
\usepackage[dvipdfmx,hmargin=16mm,vmargin=24mm]{geometry}
\fi
\usepackage{amsmath,amssymb}
\usepackage{enumitem}

\setlength{\parindent}{0pt}
\setlength{\parskip}{0.35em}
\setlist[enumerate]{leftmargin=2.2em,itemsep=0.35em,topsep=0.35em}

\newcommand{\problem}[2][10]{\par\vspace{0.85em}\noindent{\large 問題 #2}\quad（#1点）\quad\ignorespaces}
\newcommand{\R}{\mathbb{R}}

\begin{document}

\ifPDFTeX
\begin{CJK}{UTF8}{min}
\fi

\begin{center}
{\Large 数学1A レポート解答}
\end{center}

\problem{1}
$f(0,0)=0$ である。まず
\[
\partial_x f(0,0)
=\lim_{h\to 0}\frac{f(h,0)-f(0,0)}{h}=0,
\qquad
\partial_y f(0,0)
=\lim_{k\to 0}\frac{f(0,k)-f(0,0)}{k}=0.
\]
したがって，もし全微分可能なら一次近似は $0$ でなければならない。

ところが $y=x$ とおくと
\[
f(x,x)=\frac{x^3}{2x^2}=\frac{x}{2}
\]
であり，
\[
\frac{|f(x,x)-0|}{\sqrt{x^2+x^2}}
=\frac{1}{2\sqrt{2}}
\]
となって $0$ に近づかない。よって $f$ は $(0,0)$ で全微分可能でない。

\problem{2}
\[
f_x(x,y)=\frac{1}{2\sqrt{1+x+2y}},
\qquad
f_y(x,y)=\frac{1}{\sqrt{1+x+2y}}.
\]
したがって
\[
f(0,0)=1,\qquad f_x(0,0)=\frac12,\qquad f_y(0,0)=1.
\]
$f$ は $(0,0)$ の近くで $C^1$ 級なので，
\[
f(h,k)=1+\frac12h+k+o\left(\sqrt{h^2+k^2}\right)
\qquad ((h,k)\to(0,0)).
\]
よって
\[
A=\frac12,\qquad B=1.
\]

\problem{3}
$F(x,y)=x^2+xy+y^2$ とおくと
\[
F_x=2x+y,\qquad F_y=x+2y.
\]
点 $(1,-1)$ では
\[
F(1,-1)=1,\qquad F_x(1,-1)=1,\qquad F_y(1,-1)=-1.
\]
接平面は
\[
z-1=(x-1)-(y+1)
\]
である。すなわち
\[
z=x-y-1.
\]

\problem{4}
ランダウのオー記号は，関数の増大や誤差の大きさを比較する記号である。
例えば $f(x)=O(g(x))$ は，ある定数 $C>0$ により
$|f(x)|\le C|g(x)|$ と評価できることを表す。
$o(g(x))$ は $f(x)/g(x)\to 0$ を意味し，より小さい量であることを表す。

\problem{5}
$s=h+2k$ とおくと，
\[
\log(1+s)=s-\frac{s^2}{2}+o(s^2)
\qquad (s\to 0).
\]
ここで $s=O(\sqrt{h^2+k^2})$ であるから，
\[
\log(1+h+2k)
=h+2k-\frac12(h+2k)^2+o(h^2+k^2).
\]
整理すると
\[
f(h,k)
=h+2k-\frac12h^2-2hk-2k^2+o(h^2+k^2).
\]
よって
\[
A=1,\qquad B=2,\qquad C=-\frac12,\qquad D=-2,\qquad E=-2.
\]

\problem{6}
\[
f_x=3x^2-3y,\qquad f_y=3y^2-3x.
\]
停留点は
\[
y=x^2,\qquad x=y^2
\]
を満たす点である。これより
\[
x=x^4
\]
となるので $x=0,1$。したがって停留点は
\[
(0,0),\qquad (1,1)
\]
である。

ヘッセ行列は
\[
H=
\begin{pmatrix}
6x & -3\\
-3 & 6y
\end{pmatrix}.
\]
$(0,0)$ では
\[
\det H=-9<0
\]
なので鞍点である。$(1,1)$ では
\[
\det H=27>0,\qquad f_{xx}(1,1)=6>0
\]
なので極小点である。極大点は存在しない。

\problem[20]{7}
\begin{enumerate}[label=(\alph*)]
\item
\[
\nabla T(x,y)=(-2x,-4y)
\]
であるから，
\[
\nabla T(1,2)=(-2,-8).
\]

\item
$T(1,2)=100-1-8=91$ である。等温線 $T(x,y)=91$ の接線は
\[
\nabla T(1,2)\cdot ((x,y)-(1,2))=0
\]
より
\[
-2(x-1)-8(y-2)=0.
\]
したがって
\[
x+4y-9=0.
\]
\end{enumerate}

\problem[20]{8}
\begin{enumerate}[label=(\alph*)]
\item
\[
(g_a)_x=4x^3-4ay,\qquad (g_a)_y=4y^3-4ax.
\]
停留点は
\[
x^3=ay,\qquad y^3=ax
\]
を満たす点である。

$a=0$ のとき，停留点は
\[
(0,0)
\]
のみである。

$a>0$ のとき，停留点は
\[
(0,0),\qquad (\sqrt a,\sqrt a),\qquad (-\sqrt a,-\sqrt a)
\]
である。

$a<0$ のとき，停留点は
\[
(0,0),\qquad (\sqrt{-a},-\sqrt{-a}),\qquad (-\sqrt{-a},\sqrt{-a})
\]
である。

\item
ヘッセ行列は
\[
H=
\begin{pmatrix}
12x^2 & -4a\\
-4a & 12y^2
\end{pmatrix}.
\]

$a=0$ のとき，
\[
g_0(x,y)=x^4+y^4\ge 0
\]
なので，$(0,0)$ は極小点である。

$a\ne 0$ のとき，$(0,0)$ では
\[
\det H=-16a^2<0
\]
なので鞍点である。

非零の停留点では $x^2=y^2=|a|$ であるから，
\[
\det H=144a^2-16a^2=128a^2>0,\qquad
g_{xx}=12|a|>0.
\]
したがって，非零の停留点はいずれも極小点である。極大点は存在しない。
\end{enumerate}

\ifPDFTeX
\end{CJK}
\fi

\end{document}
